\thispagestyle{empty}
\vspace*{\fill}
{\centering
{\sffamily\bfseries\Large \booktitle\par}
\vspace{.35em}
{\sffamily\large \booksubtitle\par}
\vspace{.2em}
{\sffamily\small \bookdescriptor\par}
\vspace{1.5em}
{\sffamily \bookauthor\par}
}
\vspace*{\fill}
\clearpage

\thispagestyle{empty}
\begin{scopebox}
\textbf{Scope and safety.} This book begins with personal observations and develops a speculative systems framework around them. It is not medical advice, diagnosis, or proof that contemplative practice, exercise, fasting, or relaxation can regenerate a particular damaged organ. It should not be used to delay medical care or to justify extreme fasting, overtraining, forceful neck or jaw manipulation, or discontinuation of prescribed treatment.
\end{scopebox}

\vfill
{\small
Copyright \textcopyright\ 2026 Jorge Simão. All rights reserved.\\
\bookversion.\\
Manuscript developed with editorial dialogue involving Claude by Anthropic and ChatGPT by OpenAI. The author remains responsible for all claims and interpretations.
}
\clearpage

\chapter*{Preface: The Observations Came First}
\addcontentsline{toc}{chapter}{Preface: The Observations Came First}
I did not invent a theory and then search for evidence that would make it true. I practiced, observed, and only much later began to search for a language capable of describing what seemed to be happening.

For more than twenty years, Tai Chi, Qigong, meditation, stillness, movement, and deliberate relaxation gradually changed my relationship with my body. The changes were not limited to feeling calmer. I observed recurring changes in posture, breathing, balance, muscular guarding, jaw and neck organization, movement, and the relation between conscious intention and spontaneous physical adjustment. Some events were subtle. Others were audible or macroscopic. Most importantly, they appeared to continue over a timescale far longer than the brief adaptation window usually associated with relaxation.

The everyday observation behind this book is simple:

\begin{principlebox}[The observation]
When unnecessary control is progressively released, the body appears to continue reorganizing and repairing its structure.
\end{principlebox}

That sentence is an observation and an interpretation, not yet an established biological conclusion. The precise tissues, mechanisms, magnitude, and generality remain open questions. But the persistence of the observation created a conceptual problem. What kind of system can repeatedly reconstruct itself while its material components change? Why should conscious effort sometimes obstruct rather than improve organization? How can a body move toward a more coherent shape without a central engineer specifying that shape? Why does adaptation appear to occur at the level of identity, behaviour, nervous regulation, mechanics, and cells at the same time?

No single discipline seemed sufficient. Biomechanics explained loading. Neuroscience explained guarding and motor control. Developmental biology explained robust form. Geroscience explained accumulated damage. Complexity science explained attractors and feedback. Philosophy of biology supplied downward causation and biological relativity. Contemplative traditions supplied practices for suspending habitual control. The framework in this book emerged by assembling those partial explanations around the observations.

It is therefore not presented as \emph{my theory} in the proprietary sense. It is a provisional conceptual framework built from a long history of ideas, contemporary science, criticism from human and artificial interlocutors, and one unusually long personal experiment.

The purpose of this book is not to demand belief. It is to make the framework explicit enough to criticize, measure, refine, or reject.
\cleardoublepage
